\paragraph{Selection bias in model selection and test-set reuse.}
That the best of many noisy estimates is optimistic is classical in machine learning: \citet{cawley2010overfitting} show that optimising a model-selection criterion biases the performance estimate reported for the selected model, and argue for nested validation. For adaptive reuse of a holdout, \citet{dwork2015preserving} give a differential-privacy-based reusable holdout, \citet{blum2015ladder} give the Ladder, a leaderboard mechanism that reports only significant improvements, and \citet{feldman2019advantages} show that the worst-case bias of $k$ adaptively chosen classifiers shrinks as the number of classes grows, which bears on a 10-option benchmark. \citet{roelofs2019meta} measure the gap between a repeatedly queried holdout and a once-used test set empirically on Kaggle competitions. We study the same mechanism for LLM prompt and pipeline search and use these mechanisms as candidate reporting rules or as the comparison for them.

\paragraph{Inference after selection.}
\citet{andrews2024winners} study statistical inference on the selected (winning) estimate; \citet{berk2013valid} give post-selection intervals that are valid simultaneously over all possible selections; \citet{chernozhukov2013gaussian} give Gaussian approximations and a multiplier bootstrap for maxima of high-dimensional sums, which gives a simultaneous bound over $K$ configurations. These provide our analytic prediction of the gap and several of our reporting methods; we test them in a setting with discrete outcomes and strongly correlated configurations rather than assume their conditions.

\paragraph{Variability and selection of LLM prompts.}
Scores of an LLM vary with prompt design: \citet{sclar2024quantifying} report large accuracy spread across semantically equivalent formats and recommend reporting a range, and \citet{wang2024mmlupro} report smaller prompt sensitivity on MMLU-Pro than on MMLU. \citet{polo2024prompteval} estimate the \emph{distribution} of performance over many prompt templates from few evaluations, which is a different estimand from the score of the one configuration that is eventually used. Automatic tuning tools choose such a configuration: DSPy \citep{khattab2023dspy} compiles prompts and demonstrations for a user-supplied metric, MIPRO \citep{opsahlong2024mipro} jointly optimises instructions and demonstrations of multi-stage programs, and \citet{shi2024triple} cast prompt selection as fixed-budget best-arm identification. These works study how to \emph{find} a good configuration; we ask what the score of the configuration found on the tuning items says about unseen items, and how to report it.
